\section{Bối cảnh và lý do lựa chọn đề tài}
\label{sec:boi-canh}

\subsection{Nhu cầu tư vấn thủ tục trong hoạt động chuyên môn}
Xét trường hợp một giảng viên chuẩn bị tham dự hội thảo và đặt câu hỏi:
\textit{``Tôi muốn dùng kinh phí đề tài để tham dự hội thảo, cần chuẩn bị những
gì?''} Mục đích đã rõ, nhưng dữ kiện để hướng dẫn còn thiếu. Hội thảo diễn
ra ở đâu, vào thời điểm nào? Đề tài thuộc loại nhiệm vụ nào? Hoạt động tham
dự có phục vụ nhiệm vụ đang thực hiện hay không? Tùy điều kiện áp dụng của
từng quy định, những thông tin này có thể làm thay đổi hồ sơ cần chuẩn bị,
căn cứ tài chính hoặc đầu mối cần liên hệ. Một danh sách giấy tờ chung vì
thế chưa đủ để hướng dẫn trường hợp cụ thể.

Nhu cầu tương tự đặt ra khi đăng ký nhiệm vụ khoa học và công nghệ, đề nghị
điều chỉnh hoặc gia hạn, chuẩn bị nghiệm thu hay xin cử đi đào tạo, nghiên
cứu. Người hỏi bắt đầu từ công việc cần làm, trong khi căn cứ hướng dẫn được
tổ chức theo văn bản, loại nhiệm vụ và phạm vi phụ trách của từng đơn vị.
Để chuẩn bị đúng, họ cần xác định trường hợp của mình thuộc quy định nào,
cần đáp ứng những điều kiện gì và phải liên hệ đầu mối nào. Đây là nhu cầu
tư vấn mà đề tài tập trung giải quyết trong ba miền khoa học--công nghệ,
tài chính và nhân sự tại HUST.

\subsection{Các hình thức hỗ trợ hiện có}
Văn bản quy định, hướng dẫn và biểu mẫu cung cấp căn cứ để tìm hiểu thủ tục.
Tra cứu tài liệu giúp người hỏi xác định yêu cầu chung và thành phần hồ sơ;
trao đổi với chuyên viên giúp làm rõ cách áp dụng vào trường hợp cụ thể.
Những hình thức này hỗ trợ người dùng từ bước tìm hiểu quy định đến chuẩn
bị tài liệu trước khi thực hiện thủ tục tại đơn vị phụ trách.

Tại HUST, báo cáo của Giám đốc ngày 10/01/2025 ghi nhận việc triển khai
HUST Assistant tích hợp vào eHUST. Hình minh họa ở trang 7 cho thấy trợ lý
trả lời câu hỏi \textit{``Bài báo được tính điểm Knc như thế nào''}, một ví
dụ hỏi đáp liên quan đến hoạt động nghiên cứu~\cite{hust2025baocao}.

Ngày 16/10/2026, HUST còn tổ chức tập huấn sử dụng quy trình số và giới thiệu
công cụ AI Open Work cho lãnh đạo phụ trách mua sắm và cán bộ kế toán. Bài
giới thiệu mô tả khả năng hỏi đáp từ hồ sơ, rà soát mâu thuẫn giữa các văn
bản và đối chiếu với quy chế chi tiêu nội bộ~\cite{hust2026quytrinhso}.

Một nền tảng thương mại được dùng để tham khảo là MISA Agentwork. Tài liệu
sản phẩm mô tả khả năng nạp tài liệu nghiệp vụ, tổ chức câu hỏi thường gặp
và đồng bộ nguồn để trợ lý AI sử dụng khi trả lời~\cite{misa2026trithuc}.

Đối với các thủ tục thuộc phạm vi đề tài, nội dung tra cứu được cần gắn
với công việc, dữ kiện và điều kiện áp dụng của từng trường hợp.

\subsection{Vấn đề đặt ra đối với tư vấn theo tình huống}
Trở lại ví dụ hội thảo, một tài liệu có đề cập đến kinh phí tham dự chưa đủ
để kết luận người hỏi được áp dụng nội dung ấy. Câu hỏi ban đầu có thể chưa
chứa dữ kiện quyết định, còn người hỏi có thể chưa biết dữ kiện nào cần nêu.
Hệ thống tư vấn vì vậy cần xác định thông tin còn thiếu và chỉ hỏi bổ sung
khi thông tin đó ảnh hưởng đến hướng dẫn. Hỏi quá ít có thể dẫn tới chọn
sai hướng dẫn; hỏi dàn trải lại buộc người dùng cung cấp cả những thông tin
không tác động tới kết quả.

Ngay cả khi đã làm rõ tình huống, căn cứ trả lời có thể nằm ở nhiều nguồn:
quy định về nhiệm vụ nghiên cứu, điều kiện sử dụng kinh phí và nghĩa vụ liên
quan đến việc cử cán bộ đi công tác. Các nội dung tìm được phải được đối
chiếu trên cùng dữ kiện và đúng phạm vi áp dụng. Việc chọn nguồn còn phụ
thuộc vào thời điểm: ngày hỏi, ngày diễn ra hoạt động và mốc phê duyệt không
nhất thiết trùng nhau. Một nhiệm vụ đang thực hiện hoặc hồ sơ còn tồn đọng
có thể cần được xem xét theo điều khoản chuyển tiếp, thay vì tự động áp dụng
văn bản mới nhất.

Không dừng lại ở đó, ngữ cảnh tình huống còn có thể thay đổi ngay cả trong quá trình trao đổi. 
Nếu giảng viên bổ sung rằng sẽ dùng kinh phí đơn vị thay hoàn toàn cho kinh phí đề tài, những
hướng dẫn dựa trên nguồn kinh phí cũ cần được xét lại. Hệ thống phải xác
định phần nào còn áp dụng, phần nào cần thay đổi và căn cứ nào phù hợp với
nguồn mới. Việc cập nhật phải giữ nhất quán những nội dung vẫn có căn cứ,
đồng thời chỉ rõ các hướng dẫn đã được thay thế.

Để người hỏi có thể kiểm tra hướng dẫn, mỗi kết luận cần chỉ rõ dữ kiện
đang sử dụng, điều kiện được đối chiếu và đoạn nguồn hỗ trợ. Khi còn thiếu
dữ kiện hoặc căn cứ, câu trả lời phải xác định phần chưa thể kết luận và
thông tin cần bổ sung. Nhờ đó, người hỏi phân biệt được nội dung đã có cơ
sở áp dụng với nội dung còn cần làm rõ trước khi thực hiện thủ tục.

\subsection{Định hướng tiếp cận bài toán}
Từ các quan sát trên, đề tài xác định sản phẩm cần xây dựng là một ứng dụng
chatbot tiếng Việt hỗ trợ cán bộ, giảng viên tìm hiểu thủ tục trong ba miền
khoa học--công nghệ, tài chính và nhân sự tại HUST. Giao diện hội thoại phù hợp với cách
người dùng trình bày công việc bằng ngôn ngữ tự nhiên và bổ sung thông tin
trong quá trình trao đổi.

Sau quá trình hỏi đáp, người dùng cần biết quy định nào được áp dụng trong trường
hợp của mình, thông tin hoặc tài liệu nào cần bổ sung và nên thực hiện bước
nào tiếp theo. Những chỉ dẫn này hỗ trợ việc chuẩn bị thủ tục, việc quyết định sẽ phụ thuộc vào chính người hỏi hoặc các đơn vị có thẩm quyền liên quan. Hình~\ref{fig:boi-canh} mô tả vị trí
của chatbot trong bối cảnh sử dụng. Bài toán, mục tiêu và phạm vi hỗ trợ
được xác định cụ thể tại Mục~\ref{sec:bai-toan}.

\begin{figure}[H]
\centering\begin{tikzpicture}[node distance=9mm and 12mm]
\node[block,text width=31mm] (user) {Cán bộ, giảng viên\\Câu hỏi và dữ kiện};
\node[block,text width=43mm,right=of user] (bot) {Chatbot tư vấn\\theo tình huống};
\node[block,text width=30mm,right=of bot] (answer) {Hướng dẫn\\và căn cứ};
\node[block,text width=43mm,below=of bot] (source) {Nguồn KHCN, TCKH, TCNS\\và phạm vi được phép};
\node[block,text width=31mm,left=of source] (owner) {Chủ nội dung\\Xác nhận nguồn};
\draw[arrow] (user)--(bot);
\draw[arrow] (bot)--(answer);
\draw[arrow] (owner)--(source);
\draw[arrow] (source)--(bot);
% Đường hồi tiếp đi phía trên cả hai khối, đầu mũi tên chạm đỉnh khối người dùng.
\coordinate (clarify-left) at ([yshift=7mm]user.north);
\coordinate (clarify-right) at (bot.north |- clarify-left);
\draw[arrow] (bot.north) -- (clarify-right)
  -- node[midway,above,font=\footnotesize,inner sep=2pt]{Hỏi bổ sung khi cần}
  (clarify-left) -- (user.north);
\end{tikzpicture}

\caption{Bối cảnh sử dụng chatbot tư vấn theo tình huống}
\label{fig:boi-canh}
\end{figure}
